\documentclass[../../main.tex]{subfiles}
\begin{document}

\begin{flushright}
  \footnotesize\textit{【自動文字起こし・要確認】}
\end{flushright}

{\bf [解]}
\paragraph{(1)}

$1 - \left(\dfrac{5}{6}\right)^n$

\paragraph{(2)}

排反は「4が1回も出ない」かつ「2，6が0回or1回」の時で，
$2, 4, 6$ が出ない $\cdots \left(\dfrac{1}{2}\right)^n$
$4$ がでないかつ $2$ or $6$ が1回出る $\cdots {}_n C_1 \dfrac{2}{6} \left(\dfrac{1}{2}\right)^{n-1} = \dfrac{n}{3}\left(\dfrac{1}{2}\right)^{n-1}$
だからもとめるのは
\begin{align}
  q &= 1 - \left(\dfrac{1}{2}\right)^n - \dfrac{n}{3}\left(\dfrac{1}{2}\right)^{n-1} \\
  &= 1 - \left(1+\dfrac{2}{3}n\right)\left(\dfrac{1}{2}\right)^n
\end{align}
\paragraph{(3)}

$X_n$ が5でも4でもわり切れないのは，4，5が出ず，かつ2か6が1回以下出る時で，(2)と同様に
\begin{align}
  r &= \left(\dfrac{1}{3}\right)^n + {}_n C_1 \dfrac{2}{6} \left(\dfrac{1}{3}\right)^{n-1} \\
  &= (n+1)\left(\dfrac{1}{3}\right)^n
\end{align}
表から，包除定理から
\begin{align}
  p+q-p_n &= 1-r \\
  \therefore p_n &= p+q+r-1 \\
  &= 1 - \left(\dfrac{5}{6}\right)^n - \left(1+\dfrac{2}{3}n\right)\left(\dfrac{1}{2}\right)^n + (n+1)\left(\dfrac{1}{3}\right)^n
\end{align}
となり
\begin{align}
  1-p_n &= \left(\dfrac{5}{6}\right)^n + \left(1+\dfrac{2}{3}n\right)\left(\dfrac{1}{2}\right)^n \\
  &\quad - (n+1)\left(\dfrac{1}{3}\right)^n \\
  &= \left(\dfrac{5}{6}\right)^n \left\{ 1 + \left(1+\dfrac{2}{3}n\right)\left(\dfrac{3}{5}\right)^n - (n+1)\left(\dfrac{2}{5}\right)^n \right\}
\end{align}
から
\begin{align}
  \dfrac{1}{n} \log (1-p_n) = \log \dfrac{5}{6} + \dfrac{1}{n} \log \left\{ 1 + \left(1+\dfrac{2}{3}n\right)\left(\dfrac{3}{5}\right)^n - (n+1)\left(\dfrac{2}{5}\right)^n \right\} \to \log \dfrac{5}{6}
\end{align}

\end{document}
